\begin{textsample}{POS Dim 1 – 2015\_c4 – Score -3.00 – 2015\_c4/enade\_2015\_tec\_mark\_6b65a1aace.txt}  \label{ex:f1_pos_333}
CAPÍTULO 4 PERCEPÇÃO DA PROVA As análises feitas neste capítulo tratam das percepções dos concluintes da Área de Tecnologia em Marketing sobre a prova aplicada no Enade/2015.
Essas percepções foram mensuradas por meio de nove questões que avaliaram desde o grau de dificuldade da prova até o tempo gasto para concluí -la.
As percepções sobre a prova foram relacionadas com o desempenho dos estudantes e com a Grande Região de funcionamento do curso.
O questionário de percepção da prova encontra -se ao final do Anexo VII, que traz a reprodução da prova.
O desempenho dos estudantes foi classificado em quatro quartos.
Para tanto, esse desempenho foi ordenado de forma ascendente.
O percentil 25, P25, também conhecido como primeiro quartil, é a nota de desempenho que deixa um quarto ( 25 \% ) dos valores observados abaixo e três quartos acima.
A Figura 1 apresenta uma ilustração deste conceito.
O quarto inferior de desempenho é composto pelas notas abaixo do primeiro quartil.
Já o percentil 75, P75, também conhecido como terceiro quartil, é o valor para o qual há três quartos ( 75 \% ) dos dados abaixo e um quarto acima dele.
O quarto superior de desempenho é composto pelas notas iguais ou acima do terceiro quartil.
O percentil 50, P50, também conhecido como mediana, é o valor que divide as notas em dois conjuntos de igual tamanho.
O segundo quarto inclui valores entre o primeiro quartil ( P25 ) e a mediana.
O terceiro quarto contém os valores entre a mediana ( P50 ) e o terceiro quartil ( P75 ).
Vale ressaltar que percentis, quartis e medianas são pontos que não obrigatoriamente pertencem ao conjunto original de dados, ao passo que os quartos são subconjuntos dos dados originais.
A seguir, serão apresentados gráficos com resultados selecionados, relativos às nove questões avaliadas por grupos de estudantes.
Os gráficos apresentam nas barras o percentual de alunos que assinalaram uma das opções ou a soma das porcentagens daqueles que assinalaram duas ( ou três ) delas.
Por exemplo, para as questões 1 e 2, os gráficos apresentam a porcentagem total de participantes que assinalaram as opções ( D ) Difícil e ( E ) Muito difícil.
Em cada barra foram assinalados também os extremos do intervalo de confiança de 95 \% como linhas verticais unidas por uma linha horizontal na forma da letra H maiúscula.
As Tabelas no Anexo II apresentam os valores absolutos e a distribuição percentual17 das alternativas válidas das nove questões, segundo o mesmo recorte de desempenho dos alunos e Grande Região de funcionamento do curso. 4.1 GRAU DE DIFICULDADE DA PROVA 4.1.1 Componente de Formação Geral Ao avaliarem “ Qual o grau de dificuldade desta prova na parte de Formação Geral? ” ( Questão 1 ), 24,2 \% do grupo de inscritos e presentes optaram pelas alternativas Difícil ou Muito difícil.
Entretanto, para mais da metade dos estudantes ( 64,8 \% ), o Componente de Formação Geral da prova foi considerado com grau de dificuldade Médio ( Gráfico 4.1, Gráfico 4.2 e, no Anexo II, a Tabela II.1 ).
O percentual de estudantes que consideraram a prova como Difícil ou Muito difícil foi maior na região Norte, onde a proporção foi de 33,3 \%, enquanto a de menor incidência foi a Nordeste, com 20,2 \%.
No Gráfico 4.1, é possível observar diferença estatisticamente significativa apenas entre as regiões Nordeste ( 20,2 \% ) e Centro-Oeste ( 28,6 \% ).
Nas Grandes Regiões, a proporção de presentes à prova que consideraram o Componente de Formação Geral como sendo de grau de dificuldade Médio esteve entre 52,8 \% na região Norte e 65,6 \% na região Nordeste. 17 Cumpre lembrar uma das convenções para tabelas numéricas ( pág. iii ) sobre a possibilidade de a soma das partes não resultar em 100 \% por questões de arredondamento.
O percentual de alunos que consideraram a prova Difícil ou Muito difícil foi decrescente em função dos quartos de desempenho : 30,8 \% no primeiro quarto e 17,0 \% no quarto, grupo de melhor desempenho na prova.
Nos quartos de desempenho intermediários, a proporção de alunos que consideraram a prova Difícil ou Muito difícil foi de 25,9 \% no segundo quarto e 23,1 \% no terceiro quarto.
Observam -se diferenças estatisticamente significativas entre o quarto de melhor desempenho e os demais, assim como entre o quarto inferior e os demais.
A proporção de presentes que consideraram o componente de Formação Geral como sendo de grau de dificuldade Médio variou de 59,7 \% ( quarto inferior ) a 67,5 \% ( terceiro quarto ).
No quarto superior a proporção foi 66,7 \%. 4.1.2 Componente de Conhecimento Específico Ao responderem à Questão 2 – “ Qual o grau de dificuldade desta prova na parte de Componente Específico? ” – 27,5 \% do grupo de estudantes classificaram -na como Difícil ou Muito difícil.
O Componente de Conhecimento Específico da prova foi considerado com grau de dificuldade Médio por 63,6 \% dos alunos ( Gráfico 4.3, Gráfico 4.4, e, no Anexo II, a Tabela II.2 ).
A análise das respostas dos estudantes quanto ao grau de dificuldade do Componente de Conhecimento Específico da prova, agregado por Grande Região, teve maior proporção de alunos que a avaliaram como Difícil ou Muito difícil na região Norte ( 38,9 \% ) e menor proporção na região Nordeste ( 19,0 \% ), havendo diferenças estatisticamente significativas entre essa região e as demais, assim como entre a região Sudeste ( 26,7 \% ) e a Centro-Oeste ( 34,2 \% ).
O percentual de alunos que classificaram o grau de dificuldade como Médio, no Componente de Conhecimento Específico, variou de 55,6 \% a 68,4 \%, para as regiões Norte e Nordeste, respectivamente.
Considerando -se a avaliação da dificuldade das questões do Componente de Conhecimento Específico da prova, de acordo com o desempenho dos estudantes, observa -se que há diferença estatisticamente significativa do resultado do quarto de pior desempenho ( 33,4 \% ) em relação aos demais, bem como entre o terceiro quarto ( 27,5 \% ) e o quarto superior ( 22,4 \% ).
A proporção dos que classificaram a parte específica como Difícil ou Muito difícil variou de 22,4 \% ( último quarto ) a 33,4 \% ( primeiro quarto ).
A alternativa modal para a Questão 2 foi o grau Médio, com 58,3 \% do quarto inferior e 66,5 \% do superior optando por essa resposta. 4.2 EXTENSÃO DA PROVA EM RELAÇÃO AO TEMPO TOTAL Indagados quanto à extensão da prova, em relação ao tempo total oferecido para a sua resolução ( Questão 3 ), os estudantes apontaram, com maior incidência, a alternativa que considerava a extensão adequada, para todas as agregações consideradas ( Gráfico 4.5, Gráfico 4.6, e, no Anexo II, a Tabela II.3 ).
O percentual de alunos que responderam ser a extensão da prova adequada foi de 57,8 \%.
Já 38,2 \% dos inscritos presentes consideraram que a prova foi longa ou muito longa e menos do que 4 \% a avaliaram como curta ou muito curta.
Entre as Grandes Regiões a proporção daqueles que avaliaram a prova como longa ou muito longa em relação ao tempo total destinado à sua resolução variou de 33,7 \% na região Sul até 43,1 \% na região Norte, conforme mostra o Gráfico 4.5.
A diferença entre as regiões Sudeste ( 40,2 \% ) e Sul é estatisticamente significativa.
O percentual de alunos que considerou a extensão da prova adequada variou de 54,2 \% ( Norte ) a 62,5 \% ( Sul ).
Considerando -se o desempenho dos alunos, nota -se ainda que 54,6 \% consideraram a extensão da prova adequada no quarto de desempenho inferior e 61,2 \% no de melhor desempenho ( quarto superior ).
Nos quartos intermediários, esta proporção foi 55,8 \% no segundo quarto e 59,6 \% no terceiro.
No Gráfico 4.6, que mostra as proporções de estudantes que consideraram a prova longa ou muito longa em função dos quartos de desempenho, verifica -se uma variação entre 34,9 \% ( quarto superior ) e 40,9 \% ( segundo quarto ).
Observa -se diferença estatisticamente significativa entre as proporções de escolha dessa resposta no quarto superior em relação ao primeiro e ao segundo quartos. 4.3 COMPREENSÃO DOS ENUNCIADOS DAS QUESTÕES 4.3.1 Componente de Formação Geral Com relação aos enunciados das questões do Componente de Formação Geral ( Questão 4 ), as opiniões foram positivas, já que 77,3 \% dos alunos avaliados consideraram todos ou a maioria dos enunciados das questões “ claros e objetivos ” ( Gráfico 4.7, Gráfico 4.8, e, no Anexo II, a Tabela II.4 ).
Na análise regional, a percentagem de estudantes que avaliaram que todos ou a maioria dos enunciados das questões do Componente de Formação Geral estavam “ claros e objetivos ” variou de 75,3 \% na região Sudeste a 81,7 \% na região Sul, sendo essa a única diferença estatisticamente significativa.
A análise das percepções dos estudantes sobre a clareza e objetividade dos enunciados permite afirmar que todos, ou a maioria dos enunciados de questões relativas ao Componente de Formação Geral, foram considerados “ claros e objetivos ” para a maior parte dos respondentes ( maior ou igual a 75,3 \% em todas as regiões e maior ou igual a 68,6 \% para todos os quartos de desempenho ).
Segundo o desempenho, observa -se que a proporção dos que emitiram essa opinião cresce conforme o desempenho aumenta.
No quarto superior, a clareza e objetividade de todos ou da maioria dos enunciados das questões foi percebida por 85,4 \% dos alunos, observando -se diferenças estatisticamente significativas desse para os demais quartos, e também do quarto de desempenho inferior ( 68,6 \% ) em relação ao segundo ( 75,5 \% ) e ao terceiro ( 79,7 \% ) quartos. 4.3.2 Componente de Conhecimento Específico Com relação aos enunciados das questões do Componente de Conhecimento Específico da prova, para 77,7 \% dos estudantes avaliados da Área de Tecnologia em Marketing, a clareza e a objetividade ( Questão 5 ) estavam presentes em todas ou na maioria das questões ( Gráfico 4.9, Gráfico 4.10, e no Anexo II, a Tabela II.5 ).
A maioria dos estudantes de todas as Grandes Regiões brasileiras considerou “ claros e objetivos ” todos ou a maioria dos enunciados das questões do Componente de Conhecimento Específico da prova, percentual sempre maior ou igual a 74,4 \%.
A diferença entre a maior e a menor proporção, 81,7 \% ( Sul ) e 74,4 \% ( Centro-Oeste ), é estatisticamente significativa, assim como a diferença entre as regiões Sul e Sudeste ( 76,8 \% ).
A proporção de estudantes que consideraram os enunciados das questões “ claros e objetivos ” apresenta uma tendência crescente em relação ao aumento de desempenho : mais elevada no quarto superior ( 84,1 \% ) se comparada ao quarto inferior de desempenho ( 67,9 \% ).
As diferenças do primeiro para os demais quartos e do segundo para o quarto superior são estatisticamente significativas. 4.4 SUFICIÊNCIA DAS INFORMAÇÕES/INSTRUÇÕES FORNECIDAS Ao avaliarem as informações/instruções fornecidas para a resolução das questões ( Questão 6 ), 87,7 \% dos respondentes da Área de Tecnologia em Marketing de todo o Brasil afirmaram que estas eram até excessivas ou “ suficientes ” em todas ou na maioria das questões ( Gráfico 4.11, Gráfico 4.12, e, no Anexo II, a Tabela II.6 ).
Quanto à distribuição de respondentes pelas Grandes Regiões observa -se que a proporção de estudantes que consideraram as informações/instruções fornecidas até excessivas ou “ suficientes ” em todas ou na maioria das questões foi sempre igual ou superior a 83,3 \% ( região Norte ), chegando a 90,9 \% na região Sul.
Observam -se diferenças estatisticamente significativas entre as proporções da região Sul e as das regiões Sudeste ( 86,8 \% ) e Centro-Oeste ( 84,8 \% ).
Levando -se em conta o desempenho dos participantes, notam -se diferenças estatisticamente significativas entre as opiniões de estudantes do quarto inferior e superior de desempenho, bem como deles em relação aos quartos intermediários, como mostra o Gráfico 4.12.
O percentual de participantes que avaliaram “ as informações/instruções ” como até excessivas ou “ suficientes ” em todas ou na maioria das questões foi mais elevado no quarto superior ( 93,7 \% ), percentual superior à média nacional ( 87,7 \% ).
Já no quarto inferior, a suficiência das informações/instruções declarada como até excessiva, em todas ou na maioria das questões foi percebida por 80,3 \% dos respondentes.
Estas proporções são crescentes com os quartos de desempenho. 4.5 DIFICULDADE ENCONTRADA AO RESPONDER À PROVA Questionados sobre as dificuldades com as quais se depararam ao responder à prova ( Questão 7 ), 7,7 \% dos estudantes apontaram o Desconhecimento do conteúdo.
Para 43,3 \%, a Forma diferente de abordagem do conteúdo foi indicada como dificuldade.
Já a Falta de motivação para fazer a prova foi a dificuldade apontada por 11,5 \% dos respondentes.
Considerando -se todo o Brasil, 32,7 \% dos respondentes afirmaram que não tiveram qualquer tipo de dificuldade para responder à prova ( Tabela II.7 no Anexo II ).
Os Gráficos 4.13 e 4.14 apresentam os percentuais de estudantes que apontaram o Desconhecimento do conteúdo como dificuldade percebida ao responder à prova.
Na análise por Grandes Regiões, o percentual de inscritos e presentes que apontaram o Desconhecimento do conteúdo como dificuldade ao responder à prova não superou 11 \%.
Os percentuais variaram de 7,0 \% nas regiões Sudeste e Norte a 10,8 \% na Centro-Oeste, sendo a diferença entre Centro-Oeste e Sudeste a única significativa estatisticamente.
A Forma diferente de abordagem do conteúdo foi a escolha modal dos estudantes, com percentuais que variaram de 39,2 \% ( região Nordeste ) a 54,9 \% ( Norte ).
O percentual de alunos que citaram a Falta de motivação como dificuldade variou de 8,7 \% ( região Centro-Oeste ) a 12,2 \% ( região Nordeste ).
Os que declararam não ter qualquer tipo de dificuldade para responder à prova variaram de 23,9 \% na região Norte a 34,6 \% na Sul.
Com relação aos quartos de desempenho, o Desconhecimento do conteúdo foi a opção escolhida por 5,4 \% dos estudantes do quarto superior e 10,5 \% do quarto inferior.
Há diferença estatisticamente significativa do quarto inferior em relação ao terceiro e ao quarto.
A alternativa modal para os alunos, quando agregados pelos quartos de desempenho, foi que a dificuldade encontrada foi causada pela Forma diferente de abordagem do conteúdo : 42,9 \% no quarto inferior e 40,3 \% do quarto superior assim o responderam.
A indicação de Falta de motivação como dificuldade variou de 9,7 \% ( quarto superior ) a 14,6 \% ( quarto inferior ).
A proporção dos que escolheram a opção Não tive qualquer tipo de dificuldade para responder à prova variou de 26,6 \% ( quarto inferior ) a 39,5 \% ( quarto superior ). 4.6 CONTEÚDOS DAS QUESTÕES OBJETIVAS DA PROVA Ao analisarem os conteúdos das questões objetivas da prova ( Questão 8 ), um percentual muito pequeno dos estudantes avaliados, apenas 3,0 \%, afirmou que não estudou ainda a maioria desses conteúdos ( Gráficos 4.15, Gráfico 4.16, e a Tabelas II.8 no Anexo II ).
A maioria ( 77,2 \% ) afirmou ter estudado e aprendido muitos ou todos os conteúdos.
Na análise por Grande Região, a proporção de respondentes que escolheram a opção não estudou ainda a maioria desses conteúdos foi pequena.
Observa -se que nas regiões Norte ( 4,2 \% ), Nordeste ( 6,1 \% ) e Centro-Oeste ( 3,6 \% ), apesar de pequenas, as proporções foram maiores do que a média nacional ( 3,0 \% ).
Observa -se diferença estatisticamente significativa entre a região Nordeste em relação à Sul e à Sudeste.
Em todas as regiões, a maioria dos presentes afirmou ter estudado e aprendido muitos ou todos os conteúdos, com proporções variando entre 62,0 \% na região Norte e 78,7 \% na Sul.
Considerando -se separadamente as opiniões de estudantes dos quatro quartos de desempenho, observa -se que, no quarto inferior, 6,4 \% ofereceram como resposta que não estudou ainda a maioria desses conteúdos, sendo 0,9 \% os do quarto superior com a mesma resposta.
As diferenças entre os alunos que optaram por esse motivo de dificuldade no quarto superior em relação ao primeiro e ao segundo quartos, e do primeiro em relação aos quartos intermediários são estatisticamente significativas.
Tendo em conta o quarto superior, 90,4 \% dos alunos afirmaram ter estudado e aprendido muitos ou todos os conteúdos.
No outro extremo, no primeiro quarto, 61,2 \% optaram pelas mesmas categorias. 4.7 TEMPO GASTO PARA CONCLUIR A PROVA Ao responderem sobre o tempo de conclusão da prova ( Questão 9 ), mais da metade dos estudantes ( 70,4 \% ) afirmou ter gasto Entre duas e quatro horas ( Gráfico 4.17, Gráfico 4.18 e, no Anexo II, a Tabela II.9 ).
Considerando -se as cinco Grandes Regiões brasileiras, o percentual dos que utilizaram Entre duas e quatro horas para finalizar a prova apenas na região Sul ( 65,5 \% ) foi inferior ao percentual nacional.
Nas demais Grandes Regiões, o percentual de alunos que dispensaram Entre duas e quatro horas para concluir a prova, foi igual ou superior a 71 \%, como mostra o Gráfico 4.17.
Observam -se diferenças estatisticamente significativas da região Sul em relação às demais regiões.
Uma vez sendo analisadas as alternativas escolhidas pelos estudantes que se situam nos diferentes quartos de desempenho, observa -se uma tendência crescente : uma maior proporção de participantes no quarto superior ( 77,2 \% ) declarou ter gasto Entre duas e quatro horas “ para concluir a prova ” quando comparadas com os dos quartos inferiores.
As diferenças entre o primeiro ( 60,6 \% ) e o demais, bem como entre o segundo quarto ( 69,4 \% ) quando comparado ao terceiro ( 74,2 \% ) e ao quarto superior são estatisticamente significativas.

% matched lemmas: 
\end{textsample}
